\documentclass{beamer}

% Theme selection
\usetheme{Berlin}

% Essential packages
\usepackage{graphicx} % For images
\usepackage{amsmath}  % For advanced math formatting

% Presentation Metadata
\title{Introduction to Chaos}
\subtitle{Course: Complex Systems Modeling}
\author{Robert Flassig}
\institute{Brandenburg University of Applied Sciences}
\date{\today}

\begin{document}

% ---------------------------------------------------------
% Title and Outline
% ---------------------------------------------------------
\begin{frame}
    \titlepage
\end{frame}

\begin{frame}
    \frametitle{Outline}
    \tableofcontents
\end{frame}

% ---------------------------------------------------------
% Section 1: Introduction to Chaos
% ---------------------------------------------------------
\section{What is Chaos?}

\begin{frame}
    \frametitle{From Bifurcation to Chaos}
    \begin{itemize}
        \item \textbf{Recap:} We explored bifurcations (Saddle-Node, Pitchfork) where parameter changes alter stability.
        \item \textbf{Question:} What happens if we increase the parameter beyond these points?
        \item \textbf{Result:} In nonlinear systems, this often leads to \textbf{Chaos}.
        \item \textbf{Definition of Chaos:}
        \begin{itemize}
            \item Long-term behavior that never settles into a static or periodic trajectory.
            \item Deterministic (no random inputs), yet appears random.
            \item Exhibits sensitive dependence on initial conditions.
        \end{itemize}
    \end{itemize}
\end{frame}

% ---------------------------------------------------------
% Section 2: Discrete-Time Models
% ---------------------------------------------------------
\section{Chaos in Discrete-Time Models}

\begin{frame}
    \frametitle{Discrete vs. Continuous}
    \begin{itemize}
        \item \textbf{Discrete-Time ($x_t = F(x_{t-1})$):}
        \begin{itemize}
            \item Chaos is possible even in \textbf{1D} systems.
            \item Example: Logistic Map.
        \end{itemize}
        \item \textbf{Continuous-Time ($\dot{x} = f(x)$):}
        \begin{itemize}
            \item Chaos requires at least \textbf{3 dimensions} (Poincaré-Bendixson theorem).
            \item Example: Lorenz Equations.
        \end{itemize}
    \end{itemize}
\end{frame}

\begin{frame}
    \frametitle{The Logistic Map}
    A canonical model of population dynamics:
    \[
    x_t = r x_{t-1} (1 - x_{t-1})
    \]
    \begin{itemize}
        \item $x_t \in [0,1]$: Population scale.
        \item $r$: Growth rate parameter ($0 < r < 4$).
    \end{itemize}
    
    \vspace{0.5cm}
    \textbf{Behavior as $r$ increases:}
    \begin{enumerate}
        \item Stable Equilibrium (Fixed Point).
        \item Period-Doubling Bifurcation (Oscillation between 2 states).
        \item Period-Doubling Cascade ($2 \to 4 \to 8 \dots$).
        \item \textbf{Chaos} (Aperiodic behavior).
    \end{enumerate}
\end{frame}

\begin{frame}
    \frametitle{Bifurcation Diagram}
    \begin{center}
        \fbox{\parbox{0.6\textwidth}{\centering \vspace{0.5cm} \textit{Bifurcation diagram of the logistic map} \vspace{0.5cm}}}
    \end{center}
    
    \begin{itemize}
        \item \textbf{X-axis:} Parameter $r$.
        \item \textbf{Y-axis:} Asymptotic states of $x$.
        \item Shows the transition from order (single lines) to chaos (dense regions).
    \end{itemize}
\end{frame}

% ---------------------------------------------------------
% Section 3: Characteristics
% ---------------------------------------------------------
\section{Characteristics of Chaos}

\begin{frame}
    \frametitle{1. Sensitivity to Initial Conditions}
    \textbf{The "Butterfly Effect"}
    \begin{itemize}
        \item Two trajectories starting arbitrarily close together diverge exponentially fast.
        \item Divergence formula:
        \[
        |\delta(t)| \approx |\delta_0| e^{\lambda t}
        \]
        \item Implication: Long-term prediction is impossible, despite the system being deterministic.
    \end{itemize}
\end{frame}

\begin{frame}
    \frametitle{2. The Lyapunov Exponent}
    \textbf{How do we measure chaos mathematically?}
    
    \begin{itemize}
        \item The \textbf{Lyapunov Exponent ($\lambda$)} measures the rate of separation.
        \item Calculated as the time average of the log-derivative:
        \[
        \lambda = \lim_{t \to \infty} \frac{1}{t} \sum_{i=0}^{t-1} \ln \left| \frac{dF}{dx}(x_i) \right|
        \]
        \item \textbf{Criterion:} $\lambda > 0$ implies Chaos.
    \end{itemize}
\end{frame}

\begin{frame}
    \frametitle{3. Stretching and Folding}
    \textbf{The geometric mechanism of chaos:}
    
    \begin{itemize}
        \item \textbf{Stretching:} Exponential divergence of nearby points ($\lambda > 0$).
        \item \textbf{Folding:} Confinement within a finite phase space requires trajectories to fold back.
        \item This process creates \textbf{fractal} structures (Strange Attractors).
    \end{itemize}
\end{frame}

% ---------------------------------------------------------
% Section 4: Continuous Systems
% ---------------------------------------------------------
\section{Chaos in Continuous Systems}

\begin{frame}
    \frametitle{The Lorenz Equations}
    Discovered by Edward Lorenz (1963) in atmospheric modeling.
    
    \[
    \begin{aligned}
    \dot{x} &= s(y - x) \\
    \dot{y} &= rx - y - xz \\
    \dot{z} &= xy - bz
    \end{aligned}
    \]
    
    \textbf{Standard Chaotic Parameters:}
    \[
    s=10, \quad r=28, \quad b=8/3
    \]
\end{frame}

\begin{frame}
    \frametitle{The Lorenz Attractor}
    \begin{center}
        \fbox{\parbox{0.6\textwidth}{\centering \vspace{0.5cm} \textit{Lorenz attractor ($s{=}10$, $r{=}28$, $b{=}8/3$)} \vspace{0.5cm}}}
    \end{center}
    
    \begin{itemize}
        \item A \textbf{Strange Attractor}.
        \item \textbf{Properties:}
        \begin{itemize}
            \item Trajectories loop infinitely but never intersect.
            \item Has a \textbf{Fractal Dimension} (approx. 2.06).
        \end{itemize}
    \end{itemize}
\end{frame}

% ---------------------------------------------------------
% Summary
% ---------------------------------------------------------
\section{Summary}

\begin{frame}
    \frametitle{Summary}
    \begin{enumerate}
        \item \textbf{Conditions for Chaos:}
        \begin{itemize}
            \item Nonlinearity is essential.
            \item Discrete systems: Possible in 1D.
            \item Continuous systems: Requires $\ge 3$D.
        \end{itemize}
        \item \textbf{Signatures:}
        \begin{itemize}
            \item Sensitivity to initial conditions ($\lambda > 0$).
            \item Aperiodic behavior.
            \item Stretching and folding mechanism.
        \end{itemize}
    \end{enumerate}
    
    \vspace{0.5cm}
    \textbf{Next Step:} We will explore complex systems with many variables.
\end{frame}

\end{document}
